\subsection{无互异约束下的重复接受}
\label{subsec:relaxed-results}

无互异约束模型首先验证了正常批次路径的存在性（\path{normal_relaxed_batch_exists}）：至少存在一条执行能够经过批次准入并到达后续接受事件。该结果确认共同生命周期可以运行，但既不要求参与方互异，也不证明目标关系成立。

核心存在性性质\path{one_send_two_accepts_exists}同样验证通过。相应可达执行在两个槽位收集同一个$(A,\oid,m)$元组，并由一个匹配发送来源支持同批两个接受事件。关键动作事实可归纳为
\begin{equation}
  \begin{aligned}
    &\mathsf{Send}(A,\oid,m)@s,\\
    &\mathsf{BatchReceive}(\mathit{bid},\mathit{rst})@b,\\
    &\mathsf{ReceiverAccept}(A,\oid,m,\mathit{bid},\mathit{rst})@r_1,\\
    &\mathsf{ReceiverAccept}(A,\oid,m,\mathit{bid},\mathit{rst})@r_2,\\
    &s<b<r_1<r_2.
  \end{aligned}
  \label{eq:relaxed-duplicate-witness}
\end{equation}
式\eqref{eq:relaxed-duplicate-witness}中的两个接受事件具有相同的条目、批次和接收方上下文坐标，而$r_1<r_2$表明二者是不同的事件发生。该性质还约束与所示完整元组匹配的\sendaction 发生点唯一为$s$，因此两个接受并非分别对应两个匹配发送。此处的唯一性仅针对所示元组，不排除执行中存在与该元组不匹配的其他发送。

图\ref{fig:duplicate-acceptance-trace}给出了该可达执行的关键事件投影。
\begin{figure}[tbp]
  \centering
  % Audited key-event projection; see FIG2_DESIGN_REPORT.md.
% Model sid is rendered as the frozen manuscript's oid.
\begin{tikzpicture}[
  >=Latex,
  flow/.style={->,line width=0.45pt},
  read/.style={->,dashed,line width=0.35pt},
  event/.style={
    draw,rounded corners=1pt,line width=0.35pt,
    align=center,text width=5.70cm,inner sep=2.2pt,
    font=\fontsize{9}{11}\selectfont
  },
  output/.style={event,text width=3.12cm,
    font=\fontsize{9}{11}\selectfont},
  slot/.style={event,minimum height=8.5mm},
  accept/.style={event,fill=black!7,minimum height=10mm},
  time/.style={font=\fontsize{9}{11}\selectfont,inner sep=0pt},
  note/.style={font=\fontsize{7.5}{9.5}\selectfont,align=center,inner sep=1pt}
]
  % A context boundary, rather than a separate participant/actor.
  \draw[line width=0.35pt,rounded corners=1pt]
    (1.08,-2.34) rectangle (7.22,-8.66);
  \node[note,fill=white] at (4.15,-2.34)
    {同一批次上下文 $(\mathit{bid},\mathit{rst})$};

  \node[event,minimum height=11mm] (f2send) at (4.15,0) {
    匹配发送来源（该完整元组唯一）\\[-1pt]
    $\mathsf{Send}(A,\mathit{oid},m)@s$\\[-1pt]
    {\fontsize{7.5}{9.5}\selectfont 规则：\textsf{SendMessage}}
  };
  \node[output,minimum height=11mm] (f2out) at (2.60,-1.50) {
    公开条目\\[-1pt]
    $E=(A,\mathit{oid},m)$\\[-1pt]
    $\mathsf{Out}(E)$
  };
  \node[output,minimum height=11mm] (f2sent) at (6.18,-1.50) {
    持久来源事实\\[-1pt]
    $!\mathsf{Sent}(A,\mathit{oid},m)$
  };

  \node[slot] (f2slot1) at (4.15,-3.02) {
    槽位1：$E=(A,\mathit{oid},m)$\\[-1pt]
    {\fontsize{7.5}{9.5}\selectfont 规则：\textsf{CollectSlot1}；输入 $\mathsf{In}(E)$}
  };
  \node[slot] (f2slot2) at (4.15,-4.12) {
    槽位2：$E=(A,\mathit{oid},m)$\\[-1pt]
    {\fontsize{7.5}{9.5}\selectfont 规则：\textsf{CollectSlot2}；再次输入 $\mathsf{In}(E)$}
  };
  \node[event,minimum height=10mm] (f2batch) at (4.15,-5.35) {
    $\mathsf{BatchReceive}(\mathit{bid},\mathit{rst})@b$\\[-1pt]
    {\fontsize{7.5}{9.5}\selectfont 规则：\textsf{AdmitRelaxedBatch}（无互异条件）}
  };
  \node[accept] (f2accept1) at (4.15,-6.74) {
    $\mathsf{ReceiverAccept}(A,\mathit{oid},m,\mathit{bid},\mathit{rst})@r_1$\\[-1pt]
    {\fontsize{7.5}{9.5}\selectfont 规则：\textsf{ProcessSlot1}}
  };
  \node[accept] (f2accept2) at (4.15,-8.05) {
    $\mathsf{ReceiverAccept}(A,\mathit{oid},m,\mathit{bid},\mathit{rst})@r_2$\\[-1pt]
    {\fontsize{7.5}{9.5}\selectfont 规则：\textsf{ProcessSlot2}}
  };

  % One rule firing exposes E and creates one persistent origin fact.
  \draw[flow] (f2send.south) -- (f2out.north);
  \draw[flow] (f2send.south) -- (f2sent.north);

  % A single public entry is selected/delivered twice by the network.
  \draw[line width=0.45pt] (f2out.west) -- (0.88,-1.50) -- (0.88,-4.12);
  \draw[flow] (0.88,-3.02) -- (f2slot1.west);
  \draw[flow] (0.88,-4.12) -- (f2slot2.west);
  \draw[flow] (f2slot1) -- (f2slot2);
  \draw[flow] (f2slot2) -- (f2batch);
  \draw[flow] (f2batch) -- (f2accept1);
  \draw[flow] (f2accept1) -- (f2accept2);

  % Both reads have the same origin (#s:c1 in both audited graphs).
  \draw[dashed,line width=0.35pt]
    (f2sent.east) -- (8.04,-1.50) -- (8.04,-8.05);
  \draw[read] (8.04,-6.74) -- (f2accept1.east);
  \draw[read] (8.04,-8.05) -- (f2accept2.east);

  % Time points mark action events; collection rules have no action fact.
  \node[note] at (0.43,0.90) {时间};
  \draw[->,line width=0.35pt] (0.43,0.68) -- (0.43,-8.63);
  \foreach \t/\y in {s/0,b/-5.35,r_1/-6.74,r_2/-8.05} {
    \fill (0.43,\y) circle (1.25pt);
    \node[time,anchor=east] at (0.30,\y) {$\t$};
  }

  \node[font=\fontsize{9}{11}\selectfont,align=center] at (4.15,-9.07)
    {$E_1=E_2=(A,\mathit{oid},m)$\qquad $s<b<r_1<r_2$\\[-1pt]
     两个接受事件不同：$r_1\neq r_2$};
  \node[note] at (4.15,-9.72)
    {实线：条目投递/执行顺序\qquad 虚线：读取同一持久来源事实};
\end{tikzpicture}

  \caption{无互异约束模型中的重复接受可达执行}
  \label{fig:duplicate-acceptance-trace}
  \par\smallskip
  \begin{minipage}{8.2cm}
    \fontsize{7.5}{10}\selectfont
    依据Tamarin模型及独立复核的可达执行重构，仅显示与所示完整元组相关的关键事件；该元组的匹配Send唯一，执行中可能存在与其不匹配的其他Send。
  \end{minipage}
\end{figure}

该重复接受源于批次组合阶段对同一已暴露条目的重复选取：同一条目占据两个槽位，两个处理步骤随后读取同一持久来源事实并分别产生接受事件。因此，重复参与方条目能够进入同一批次并产生接受事件，RQ1得到肯定回答；一个匹配发送来源也确实能够对应同批两个接受事件，RQ2得到肯定的存在性回答。

与之对应的全称性质\path{receiver_accept_injective}被反例否定。式\eqref{eq:relaxed-duplicate-witness}满足同一发送早于两个同坐标接受事件的前件，但$r_1\neq r_2$。因此，该执行既是\path{one_send_two_accepts_exists}的可达见证，也是批内来源单射性的反例轨迹，两项结果描述同一种重复形状。与此同时，\path{receiver_accept_has_send}验证通过，表明两个接受均有更早的匹配发送。该反例揭示的是具有模型内匹配发送来源的条目在两槽抽象中被重复使用，而非无来源接受或满足K-Waay原有互异条件的完整协议攻击。

\subsection{消息互异约束的替代性分析}
\label{subsec:message-results}

消息互异约束模型中的\path{accepted_batch_has_distinct_messages}验证通过：同一$(\mathit{bid},\mathit{rst})$中的两个不同\receiveraccept 事件具有不同消息分量。该全称性质确认了模型直接施加的批内消息关系，但不涉及参与方分量。

该模型的发送来源对应性和批内来源单射性也均验证通过。前者表明每个接受事件都有更早的完整元组匹配发送，后者排除一个匹配发送来源在同一批次和接收方上下文中对应两个不同接受事件。三项性质分别记录消息比较、匹配发送存在性和同一来源所对应的接受事件数，均不比较两个条目的参与方分量。按照第\ref{subsec:configurations-properties}节的规则语义，批内消息互异性与批内来源单射性在当前模型中相互蕴含，不能作为两项独立的累积保证。

替代性结论的关键证据是\path{same_party_different_messages_batch_exists}的验证结果。相应可达执行包含同一参与方$A$的两个条目$(A,\oid_1,m_1)$和$(A,\oid_2,m_2)$，其中$\oid_1\neq\oid_2$且$m_1\neq m_2$。两个条目各有自己的匹配$\mathsf{Send}$来源，且两个发送都早于同一次批次准入；随后，两个条目通过消息级准入，并在共同$(\mathit{bid},\mathit{rst})$中依次产生接受事件。该执行没有复用匹配发送来源，重复只出现在参与方投影上。

该机器见证将$m_1\neq m_2$且$A_1=A_2$的静态关系落实到完整的共同生命周期，并确认两个接受各有匹配发送来源。因此，消息级准入约束不能推出式\eqref{eq:distinct-party-objective}规定的批内参与方互异性。批内消息互异性和批内来源单射性的全称结果仍分别对应各自公式与复核记录，但不构成相互独立的增量证据。消息相等时的拒绝分支也存在可达执行（\path{repeated_message_rejection_exists}）；该结果仅说明分支可达，不保证所有重复消息最终均被拒绝。

由此，消息互异约束模型能够维护自身断言的消息关系和来源单射关系，却不能替代参与方互异条件。RQ3的否定答案来自同一参与方不同消息批次的实际可达性，而不是其他性质的验证失败。

\subsection{参与方互异约束下的性质验证}
\label{subsec:party-results}

参与方互异约束模型在准入时直接要求$A_1\neq A_2$。目标全称性质\path{accepted_batch_has_distinct_parties}验证通过，表明同一$(\mathit{bid},\mathit{rst})$中的两个不同\receiveraccept 事件具有不同参与方投影。准入规则规定执行如何受到目标条件约束，全称性质则验证后续接受事件满足的关系；二者分别属于规则语义和机器验证结论。

存在性性质\path{distinct_party_batch_exists}也验证通过。其可达执行包含两个不同参与方，两个条目各有模型内匹配发送来源；相应发送均早于同一次批次准入，随后产生两个接受事件。该结果说明有效的不同参与方批次能够完成后续处理，目标性质并非依赖拒绝全部批次而成立。

发送来源对应性与批内来源单射性均验证通过，表明有效路径中的来源关系与目标性质相容。同一参与方拒绝分支也存在可达执行（\path{same_party_rejection_exists}）。不过，该性质没有把此前的具体发送元组绑定到最终被拒的候选元组，因而不能据此断言任意同参与方输入最终都会被拒绝。

目标对照由此提供两类证据：批内参与方互异性在全部执行轨迹上成立，有效的不同参与方批次也确实可达。在当前两槽抽象中，直接约束参与方分量能够维持目标关系，同时保留有效行为。

\subsection{综合验证结果}
\label{subsec:combined-results}

表\ref{tab:key-verification-results}汇总7项核心结果。存在性性质验证通过表示至少存在一条满足公式的执行轨迹；全称性质验证通过表示模型的所有执行轨迹均满足相应断言；全称性质被反例否定则表示存在违反断言的反例轨迹。三种结果必须按照各自量词解释。

\begin{table}[htbp]
  \centering
  \caption{关键Tamarin验证结果}
  \label{tab:key-verification-results}
  \small
  \begin{tabularx}{\linewidth}{@{}L{0.08\linewidth}L{0.40\linewidth}L{0.10\linewidth}L{0.12\linewidth}Y@{}}
    \toprule
    配置 & 验证性质 & 类型 & 结果 & 论证作用 \\
    \midrule
    放宽 & 一个精确来源支持两个同批接受\newline
      {\footnotesize\path{one_send_two_accepts_exists}} & 存在迹 & 已验证 & 重复接受见证 \\
    放宽 & 作用域内发生注入性\newline
      {\footnotesize\path{receiver_accept_injective}} & 全迹 & 被反例否定 & 对应全迹性质失败 \\
    \midrule
    消息级 & 消息区分\newline
      {\footnotesize\path{accepted_batch_has_distinct_messages}} & 全迹 & 已验证 & 维护消息坐标性质 \\
    消息级 & 作用域内发生注入性\newline
      {\footnotesize\path{receiver_accept_injective}} & 全迹 & 已验证 & 排除精确重复形状 \\
    消息级 & 同一参与方不同消息批次可达\newline
      {\footnotesize\path{same_party_different_messages_batch_exists}} & 存在迹 & 已验证 & 支撑非替代性结论 \\
    \midrule
    参与方级 & 参与方区分\newline
      {\footnotesize\path{accepted_batch_has_distinct_parties}} & 全迹 & 已验证 & 维持目标参与方关系 \\
    参与方级 & 有效不同参与方批次可达\newline
      {\footnotesize\path{distinct_party_batch_exists}} & 存在迹 & 已验证 & 排除真空成立 \\
    \bottomrule
  \end{tabularx}
  \par\smallskip
  \begin{minipage}{\linewidth}
    \footnotesize 注：存在迹对应exists-trace，全迹对应all-traces；“已验证”对应verified，“被反例否定”对应falsified。性质名下方的小号标识用于追踪对应lemma。
  \end{minipage}
\end{table}


三种模型形成连续的受控比较。无互异约束模型允许一个匹配发送来源对应同批两个接受事件，并产生批内来源单射性反例；消息互异约束模型排除这种来源复用形状，保持批内消息互异性和批内来源单射性，却仍允许同一参与方的不同消息进入同一批次；参与方互异约束模型直接保持目标关系，并保留有效批次。关键可达见证支持的结论是：消息级约束不能推出批内参与方互异性；当前模型中的批内消息互异性与批内来源单射性具有规则依赖，不能作为两项独立的累积证据。

其余7项支持结果均验证通过，包括无互异约束模型的正常路径、三个模型各自的发送来源对应性、消息互异约束和参与方互异约束模型的拒绝分支可达性，以及参与方互异约束模型的批内来源单射性。这些结果分别排除生命周期不可运行或无来源接受的解释，确认相应拒绝分支可达，并补充目标对照的来源关系。全部14项结果仍为13项验证通过、1项被反例否定，结果状态和证明步数均未改变。

后续独立复核重新执行了三个未修改模型，14项结果的状态和证明步数均与既有记录一致。相应记录构成独立复核证据，并非对历史运行日志的恢复。
% TODO(arocmag-anonymous-artifact-url): acceptance 前补入匿名复现材料URL；正文不暴露本地仓库路径。

